% =============================================================================
% Pakete und Einstellungen
% =============================================================================

% --- Grundlegende Pakete ---
\usepackage[utf8]{inputenc}
\usepackage[T1]{fontenc}
\usepackage[ngerman]{babel}
\usepackage{lmodern}
\usepackage[protrusion=true, expansion=true]{microtype}

% --- Seitenränder und Layout ---
\usepackage[a4paper, left=3cm, right=3cm, top=2.5cm, bottom=2.5cm]{geometry}
\usepackage{setspace}

% Kleiner zusaetzlicher Abstand nach jeder Abschnittsueberschrift (KOMA afterskip)
\RedeclareSectionCommand[afterskip=1.5em plus .2em]{section}

% Halber Zeilenabstand zwischen Absaetzen; letzte Absatzzeile bleibt rechts frei,
% damit Absatzwechsel im Fliesstext sichtbar sind (KOMA parskip)
\KOMAoptions{parskip=half}

% --- Kopf- und Fußzeilen ---
\usepackage[headsepline]{scrlayer-scrpage}
\pagestyle{scrheadings}
\clearpairofpagestyles
\ihead{\leftmark}
\ohead{\thepage}

% --- Mathematik ---
\usepackage{amsmath}
\usepackage{amssymb}
\usepackage{amsfonts}

% --- Grafiken und Farben ---
\usepackage{graphicx}
\usepackage[table,xcdraw]{xcolor}
\usepackage{float}
\usepackage{subcaption}

% --- Tabellen ---
\usepackage{booktabs}
\usepackage{longtable}
\usepackage{tabularx}
\usepackage{multirow}
\usepackage{makecell} % \Xhline fuer horizontale Trenner variabler Dicke

% --- Listings / Quellcode ---
\usepackage{listings}
\lstset{
    basicstyle=\ttfamily\small,
    breaklines=true,
    frame=single,
    numbers=left,
    numberstyle=\tiny,
    tabsize=4,
    captionpos=b,
    showstringspaces=false
}

% --- Verzeichnisse und Verweise ---
\usepackage[hidelinks]{hyperref}
\usepackage[nameinlink]{cleveref}
% Deutsche Bezeichner fuer cleveref (\cref / \Cref)
\crefname{figure}{Abbildung}{Abbildungen}
\Crefname{figure}{Abbildung}{Abbildungen}
\crefname{table}{Tabelle}{Tabellen}
\Crefname{table}{Tabelle}{Tabellen}
\crefname{section}{Abschnitt}{Abschnitte}
\Crefname{section}{Abschnitt}{Abschnitte}
\crefname{subsection}{Abschnitt}{Abschnitte}
\Crefname{subsection}{Abschnitt}{Abschnitte}
\crefname{chapter}{Kapitel}{Kapitel}
\Crefname{chapter}{Kapitel}{Kapitel}
\crefname{equation}{Gleichung}{Gleichungen}
\Crefname{equation}{Gleichung}{Gleichungen}
\usepackage{url}
\usepackage{bookmark}

% --- Literaturverzeichnis ---
\usepackage[backend=biber, style=authoryear, sorting=nyt, maxbibnames=99, maxcitenames=1]{biblatex}
\addbibresource{literatur.bib}

% --- Abkürzungsverzeichnis ---
\usepackage[acronym, toc, nonumberlist]{glossaries}
\makeglossaries

% --- Sonstiges ---
\usepackage{csquotes}
\usepackage{enumitem}
\usepackage{caption}
\usepackage{pdflscape}
\usepackage{pdfpages}
\usepackage{eso-pic}
\usepackage{appendix}

% --- Eigene Befehle ---
\newcommand{\todo}[1]{\textcolor{red}{\textbf{TODO:} #1}}
